\section*{Bab \ref{ch:g8:nervous-communication} --- Komunikasi Saraf}

\begin{solution}{exo:g8:nervous-communication:1}
\omterm{def:g8:nervous-communication:neuron}{Neuron} adalah \omterm{def:g5:first-look-at-cells:cell}{sel} kerja \omterm{def:g5:brain-in-command:system}{sistem sarafnya}: sebuah badan \omterm{def:g5:first-look-at-cells:cell}{sel}
dengan serat panjang yang halus --- sebagian dari \omterm{def:g4:classifying-animals:vertebrate}{tulang
belakang} sampai jari kaki. \omterm{def:g8:nervous-communication:neuron}{Saraf} adalah kabel yang membundel
ribuan serat semacam itu.
\end{solution}

\begin{solution}{exo:g8:nervous-communication:2}
Isyarat listrik yang singkat, sampai \qty{100}{m/s}.
Isyaratnya semuanya serupa; kekuatannya disandikan sebagai
lajunya --- dari derap yang lembut sampai genderang yang cepat.
\end{solution}

\begin{solution}{exo:g8:nervous-communication:3}
\omterm{def:g8:nervous-communication:neuron}{Neuron} indera: jalur laporannya, dari \omterm{def:g1:the-five-senses:sense}{alat indra} ke sumsum
dan \omterm{def:g5:brain-in-command:system}{otaknya}. \omterm{def:g8:nervous-communication:neuron}{Neuron} gerak: jalur perintahnya, dari sumsum ke
\omterm{def:g3:bones-and-muscles:muscle}{otot}. \omterm{def:g8:nervous-communication:neuron}{Neuron} penghubung: pemutusannya, di dalam \omterm{def:g5:brain-in-command:system}{otak} dan
sumsumnya.
\end{solution}

\begin{solution}{exo:g8:nervous-communication:4}
Celah renik tempat ujung serat sebuah \omterm{def:g8:nervous-communication:neuron}{neuron} bertemu \omterm{def:g5:first-look-at-cells:cell}{sel}
berikutnya. Sebuah pengantar pesan kimiawi menyeberanginya,
berlabuh, lalu --- bila cukup banyak yang tiba --- memulai
lagi isyaratnya di dalam \omterm{def:g8:nervous-communication:neuron}{neuron} berikutnya.
\end{solution}

\begin{solution}{exo:g8:nervous-communication:5}
Keterstelan (sinapsisnya menguat karena dipakai dan memudar
bila tidak --- itulah belajar); lalu lintas searah
(pelepasannya hanya terjadi di satu sisi); dan penimbangan
(tiap \omterm{def:g8:nervous-communication:neuron}{neuron} menjumlahkan ribuan masukan yang memacu dan
meredam sebelum menyalak).
\end{solution}

\begin{solution}{exo:g8:nervous-communication:6}
\omterm{def:g8:nervous-communication:neuron}{Saraf}: milidetik, satu alamat yang tepat, jalur kabel
khusus, akibatnya lenyap bersama isyaratnya. \omterm{def:g8:puberty:hormones}{Hormon}: menit
sampai tahun, seluas tubuh, siaran aliran \omterm{prop:g4:journey-of-food:absorb}{darah}, akibatnya
berjalan selama pengantar pesannya beredar.
\end{solution}

\begin{solution}{exo:g8:nervous-communication:7}
\omterm{def:g8:nervous-communication:neuron}{Neuron} indera \omterm{def:g1:the-five-senses:sense}{matanya} menggenderangkan kedudukan bolanya ke
\omterm{def:g5:brain-in-command:system}{otaknya}; \omterm{def:g8:nervous-communication:neuron}{neuron} penghubungnya menimbang arusnya di
sinapsisnya --- jalur yang dihitung itu adalah masukan yang
dijumlahkan simpang demi simpang; \omterm{def:g8:nervous-communication:neuron}{neuron} geraknya
menggenderangkan perintahnya keluar, dan lajunya menetapkan
kekuatan tiap \omterm{def:g3:bones-and-muscles:muscle}{ototnya} untuk langkah, condong dan mengatupnya
tangannya.
\end{solution}

\begin{solution}{exo:g8:nervous-communication:8}
Tiap lemparan yang dilatih menghantarkan isyaratnya
melintasi simpang yang sama, dan simpang yang diseberangi
menguat: sesudah sebulan perutean sulapnya sudah bertimbangan
siap --- yang tadinya menuntut penimbangan yang melelahkan
kini berjalan sendiri. Ketetapannya adalah \omterm{def:g8:nervous-communication:synapse}{sinapsis} yang
lebih kuat.
\end{solution}

\begin{solution}{exo:g8:nervous-communication:9}
Karena sinapsisnyalah tempat isyaratnya menjadi kimia ---
sebuah pengantar pesan yang dilepaskan, tempat berlabuh, dan
sebuah celah yang terbuka bagi molekul apa pun yang tiba.
Listrik pada sebuah kabel sukar dipalsukan; kurir kimiawi
dapat ditirukan, dihadang atau dikuras --- di situlah obat
bius, bisa dan obat-obatan sama-sama bekerja.
\end{solution}

\begin{solution}{exo:g8:nervous-communication:10}
(a) Milidetik, \omterm{def:g3:bones-and-muscles:muscle}{otot} satu \omterm{def:g1:my-body:parts}{lengan}, jalur berkabel, lenyap
seketika: \omterm{def:g8:nervous-communication:neuron}{saraf}. (b) Bertahun-tahun, seluruh tubuh, aliran
\omterm{prop:g4:journey-of-food:absorb}{darah}, berjalan berbagai musim: \omterm{def:g8:puberty:hormones}{hormon}. (c) Berantai: awal
sepersekian detik yang bersifat \omterm{def:g8:nervous-communication:neuron}{saraf}, lalu pengantar pesan
aliran \omterm{prop:g4:journey-of-food:absorb}{darahnya} menahan siaganya beberapa menit sesudah
kagetnya --- keduanya, berurutan.
\end{solution}

\begin{solution}{exo:g8:nervous-communication:11}
Perintahnya tiba di simpangnya lalu tak dapat menyeberang:
\omterm{def:g3:bones-and-muscles:muscle}{ototnya} tidak menerima apa pun. Korbannya lumpuh sedikit
demi sedikit --- sadar, merasa, memutuskan --- dengan
gelungnya terputus pada mata rantai terakhirnya,
perintah-ke-otot.
\end{solution}

\begin{solution}{exo:g8:nervous-communication:12}
Belajar: keterampilan adalah \omterm{def:g8:nervous-communication:synapse}{sinapsis} yang menguat karena
dipakai --- sebulan sang penyulap bola hidup di simpangnya.
Memutuskan: tiap \omterm{def:g8:nervous-communication:neuron}{neuron} menyalak atau tidak dengan
menjumlahkan suara simpangnya --- ``memutuskan'' pada
gelungnya adalah hitungan simpang. Kerawanan: obat bius,
bisa dan obat-obatan bekerja di tempat isyaratnya menjadi
kimia --- di celahnya. Kabelnya hanya mengangkut; simpangnya
belajar, memilih dan dapat diperdaya.
\end{solution}

\begin{solution}{pb:g8:nervous-communication:1}
\textbf{1.} Regangan uratnya menyalakan \omterm{def:g8:nervous-communication:neuron}{neuron} inderanya;
\omterm{def:g8:nervous-communication:neuron}{neuron} itu berjalan ke \emph{\omterm{def:g5:brain-in-command:system}{sumsum tulang belakang}},
bersinapsis di situ langsung pada \omterm{def:g8:nervous-communication:neuron}{neuron} geraknya, dan
genderang geraknya menendangkan \omterm{def:g3:bones-and-muscles:muscle}{otot} pahanya. Dua jenis
\omterm{def:g8:nervous-communication:neuron}{neuron}, satu simpang, setingkat sumsum --- \omterm{def:g5:brain-in-command:system}{otaknya} dipintas
karena perlindungan tak dapat menunggu pertimbangan.

\textbf{2.} Gelung tangkapannya berjalan dari \omterm{def:g1:the-five-senses:sense}{mata} ke \omterm{def:g5:brain-in-command:system}{otak}
ke \omterm{def:g1:my-body:parts}{lengan} --- bermeter-meter serat dan banyak simpang
penimbang; sentakan \omterm{def:g1:my-body:joint}{lututnya} berjalan dari urat ke sumsum ke
\omterm{def:g3:bones-and-muscles:muscle}{otot} --- beberapa sentimeter dan pada dasarnya satu \omterm{def:g8:nervous-communication:synapse}{sinapsis}.
Kabel yang lebih pendek, simpang yang lebih sedikit,
sepersepuluh waktunya.

\textbf{3.} Keputusannya dibuat di dalam sumsumnya, sebelum
veto \omterm{def:g5:brain-in-command:system}{otak} mana pun dapat tiba --- dan lalu lintas searah
sinapsisnya berarti keberatan \omterm{def:g5:brain-in-command:system}{otaknya} yang datang kemudian
tak dapat berjalan mundur menuruni jalur inderanya.
Peringatan tidak mengubah apa pun: gelungnya menutup di
bawah kehendaknya.

\textbf{4.} Rasanya adalah perjalanan kedua: laporan ketukan
yang sama juga mendaki sumsumnya ke \omterm{def:g5:brain-in-command:system}{otaknya} --- jalur yang
lebih panjang dengan lebih banyak simpang --- lalu tiba
ketika tendangannya sudah menjadi sejarah.

\textbf{5.} Penyandian frekuensinya: regangan yang lebih
keras mengirimkan genderang yang lebih cepat, dan jalur
geraknya menjawab setimpal --- kontraksi yang lebih kuat
dari laju yang lebih tinggi, bahasa satu kata di kedua
ujungnya.

\textbf{6.} Bahwa seluruh litarnya sehat: jalur inderanya,
simpang sumsumnya, jalur geraknya dan \omterm{def:g3:bones-and-muscles:muscle}{ototnya} --- dan
ketangkupannya mengesahkan kedua sisinya serupa. Satu
ketukan mengaudit empat bagian.

\textbf{7.} Bahwa suara \omterm{def:g8:nervous-communication:neuron}{neuron} penghubung dari tempat lain
di dalam sistemnya mencapai bahkan simpang gelung ini, lalu
menaikkan kesiapannya: gerak refleks yang ``tanpa \omterm{def:g5:brain-in-command:system}{otak}'' itu
tetap terangkai ke dalam keseluruhannya, dan sinapsisnya
mendengar lebih banyak daripada uratnya saja.

\textbf{8.} Sinapsisnya --- dan pengisyaratan seratnya pada
umumnya: obat biusnya menenangkan kimia di simpangnya dan
lalu lintas pesannya sekalian, tendangan dan rasanya sama-sama.

\textbf{9.} Melapor (regangan uratnya) --- memutuskan (satu
simpang sumsumnya) --- memerintah (jalur geraknya):
gelungnya dengan langkah tengahnya menyusut menjadi satu
\omterm{def:g8:nervous-communication:synapse}{sinapsis}, dan penimbangan \omterm{def:g5:brain-in-command:system}{otaknya} dipintas.

\textbf{10.} Gerak refleksnya: urat-sumsum-otot, satu
simpang, seperlima puluh detik, tak dapat ditahan, tanpa
latihan. Tangkapannya: mata-otak-lengan, simpang yang tak
terhitung, seperlima detik, sepenuhnya disengaja, dibangun
oleh latihan.

\textbf{11.} Karena penimbangan \omterm{def:g5:brain-in-command:system}{otaknya} memakan waktu, dan
sebuah luka bakar makin dalam tiap milidetik: gelung yang
melindungi jaringan dipatri pendek, setingkat sumsum, lalu
memberi tahu \omterm{def:g5:brain-in-command:system}{otaknya} sesudahnya --- terlambat merasakannya,
sudah terselamatkan.

\textbf{12.} ``Satu ketukan palu memperlihatkan bahasa
frekuensi satu katanya, kabel dua jalurnya menembus
sumsumnya, dan sebuah simpang yang memutuskan lebih cepat
daripada pemiliknya.''
\end{solution}
